\subsection{Co-occurrence Constraint and Canonical Schema}
\label{subsec:ingestion}

Genuine learned fusion is valid only where modalities co-occur on the same subject at the same
time. In the open data used here that holds for wearable physiology with CGM and meal context
(CGMacros), EEG with peripheral physiology (DEAP), and chest/wrist physiology (WESAD). No open
cohort co-registers EEG, CGM, and EHR on one subject. The serving router therefore refuses
unmatched clamp-EEG concatenated onto CGMacros people; that join is not a parent-DVXR Goal~3
object.

Every source maps onto one long-format event table before any encoder runs. The required floor is
thirteen columns (\texttt{subject\_id}, \texttt{session\_id}, \texttt{timestamp\_utc},
\texttt{source\_system}, \texttt{device}, \texttt{modality}, \texttt{channel}, \texttt{value},
\texttt{unit}, \texttt{sampling\_rate\_hz}, \texttt{quality\_flag}, \texttt{label\_name},
\texttt{label\_value}); dataset-specific extras are preserved after that floor. Windowing is
strictly no-future-leakage: features at prediction time $t$ use history $\le t$, and the target
is held out of the feature set. Splits are patient-disjoint, so temporal causality is paired with
the absence of subject leakage. This is a no-future-information property, not an interventional
causal-effect claim.

Galea and EMOTIV sessions ingested on 2026-06-08 validate the schema and plug-and-play path only:
they are single-subject, unused as training or validation labels, and contribute no reported
metric.
